\documentclass[11pt]{article}
\usepackage[margin=1in]{geometry}
\usepackage[T1]{fontenc}
\usepackage{amsmath,amssymb,amsthm,mathtools}
\usepackage[hidelinks]{hyperref}
\usepackage{enumitem}
\newtheorem{theorem}{Theorem}
\newtheorem{lemma}[theorem]{Lemma}
\newtheorem{proposition}[theorem]{Proposition}
\newtheorem{corollary}[theorem]{Corollary}
\newcommand{\Reg}{\operatorname{Reg}}
\newcommand{\TV}{\operatorname{TV}}
\newcommand{\ip}[2]{\langle #1,#2\rangle}
\title{A refined constant-regret upper bound, sparse interactions,\texorpdfstring{\\}{ }and exact minimax results for separable games}
\author{OpenAI GPT 6 Astra Pro}
\date{18 October 2026}
\begin{document}
\maketitle

\section*{Problem and claim}
\noindent\textbf{Permanent problem ID:} WS-107790.\quad
\textbf{Intake handle:} ECON-1.\\
\textbf{Problem:} Optimal dimension dependence of constant self-play regret.\\
\textbf{Claim type:} Partial result.\\
\textbf{Status of the full problem:} Unresolved.

\paragraph{Exact scope.}
Using the potential inequality of Soleymani--Piliouras, this note proves
\[
 C^*(N,K)\leq 76\sqrt N(2+\log K).
\]
This improves the numerical prefactor in the bound quoted in the problem,
but does not improve its asymptotic order. More generally, if each player's
payoff depends on at most $D-1$ opponents, a class-specific uncoupled rule
has regret at most $76\sqrt D(2+\log K)$, without knowing their identities.
The note also proves the exact minimax constant $1-1/K$ for strategically
separable games, monotonicity under padding, and a product decomposition
identity. None supplies a dimension-growing lower bound for the full class.

\paragraph{Assumptions and limitations.}
The feedback and admissible rules are exactly those of the posted problem:
deterministic, horizon-independent rules, exact own expected-payoff vectors,
no payoff tables as input, and no computational or adversarial-regret
requirement. All payoffs lie in $[0,1]$ and all logarithms are natural.
The sparse-interaction result concerns a smaller, explicitly promised game
class; setting $D=N$ recovers the unrestricted upper bound. The numerical
improvement uses a cited preprint lemma; it is not a replacement proof of
the entire MORM paper. No priority claim is made for the observations here.

\section{The exact target}
For $m=(m_1,\ldots,m_N)$, $1\leq m_i\leq K$, let
$\mathcal G_{N,m}$ and $\mathfrak A_{N,m}$ have the meanings in the problem.
For a fixed tuple of rules $\mathcal A$, put
\[
 R_i(T)=\max_{a\in A_i}\sum_{t=1}^T
 \bigl(v_i^t(a)-\ip{x_i^t}{v_i^t}\bigr),
 \qquad v_i^t(a)=u_i(a,x_{-i}^t).
\]
Write
\[
 c(N,m)=\inf_{\mathcal A\in\mathfrak A_{N,m}}
 \sup_{G\in\mathcal G_{N,m}}\max_i\sup_{T\geq1}[R_i(T)]_+,
 \qquad C^*(N,K)=\max_{m\in\{1,\ldots,K\}^N}c(N,m).
\]
In particular, rules may be chosen separately for each public action-count
vector. The bounds below respect that quantifier order.

\section{A payoff-variation inequality}
For distributions $p,q$ on a finite set define
\[
 H(p,q)^2=\sum_a(\sqrt{p_a}-\sqrt{q_a})^2,
 \qquad \TV(p,q)=\frac12\sum_a|p_a-q_a|.
\]
Suppose player $i$ has a set $S_i\subseteq\{1,\ldots,N\}\setminus\{i\}$
such that $u_i$ depends only on the actions of $i$ and $S_i$.
Set $J_i=S_i\cup\{i\}$ and define its centered vector at profile $x$ by
\[
 r_i(x)(a)=u_i(a,x_{S_i})-u_i(x_i,x_{S_i}).
\]

\begin{lemma}\label{lem:variation}
For any profiles $x,y$,
\[
 \|r_i(x)-r_i(y)\|_\infty^2
 \leq4\sum_{j\in J_i}H(x_j,y_j)^2.
\]
\end{lemma}
\begin{proof}
For any $f$ with values in $[0,1]$, decomposition into the positive and
negative parts of $p-q$ gives
$|\mathbb E_p f-\mathbb E_q f|\leq\TV(p,q)$.
Cauchy--Schwarz gives
\[
 \TV(p,q)\leq\tfrac12
 \Bigl(\sum_a(\sqrt p_a-\sqrt q_a)^2\Bigr)^{1/2}
 \Bigl(\sum_a(\sqrt p_a+\sqrt q_a)^2\Bigr)^{1/2}
 \leq H(p,q).
\]
The affinity $\rho(p,q)=\sum_a\sqrt{p_aq_a}$ belongs to $[0,1]$.
For products, the affinity is the product of the individual affinities.
The inequality $1-\prod_j\rho_j\leq\sum_j(1-\rho_j)$ therefore implies
\[
 H\Bigl(\bigotimes_{j\in J}x_j,\bigotimes_{j\in J}y_j\Bigr)^2
 \leq\sum_{j\in J}H(x_j,y_j)^2.
\]
Apply these facts separately to the two terms defining $r_i$.
For each $a$, the first term changes by at most
$\sqrt{\sum_{j\in S_i}H(x_j,y_j)^2}$; the second changes by at most
$\sqrt{\sum_{j\in J_i}H(x_j,y_j)^2}$.
Their sum is at most twice the latter square root. Squaring and maximizing
over $a$ proves the claim.
\end{proof}

\section{The refined upper bound}
For completeness, specify the exact learning rule. Put $c=2+\log K$ and
\[
 f(z)=\begin{cases}(1-z)^{-1},&z\leq0,\\1+z,&z\geq0,\end{cases}
 \qquad
 \Psi_i(U)=c\left(\sum_{a\in A_i}f(U_a/c)^{c-1}\right)^{1/(c-1)}.
\]
Set $U_i^1=0$ and $r_i^0=0$. With a fixed $0<\eta\leq1/16$, use
the MORM update
\[
 x_{ia}^t=
 \frac{\partial_a\Psi_i(\eta U_i^t)(1+4\eta r_{ia}^{t-1})}
 {\sum_b\partial_b\Psi_i(\eta U_i^t)(1+4\eta r_{ib}^{t-1})},
 \quad r_i^t=v_i^t-\ip{x_i^t}{v_i^t}\mathbf1,
 \quad U_i^{t+1}=U_i^t+r_i^t.
\]
All factors are positive: $|r_{ia}^{t-1}|\leq1$ and
$1+4\eta r_{ia}^{t-1}\geq3/4$. The first distribution is uniform.
This is MORM with a different fixed learning rate, not plain Optimistic
Hedge.

The following is the only substantive imported learning estimate.
Soleymani--Piliouras, Lemma 4.1 and Proposition 4.7~\cite{morm}, give
$\Psi_i(U)>0$, $\Psi_i(U)\geq c+\max_aU_a$, and, for every horizon,
\begin{equation}\label{eq:rvu}
 \Psi_i(\eta U_i^{T+1})+\tfrac14 V_i(T)
 \leq 3c+10\eta^2\sum_{t=1}^T\|r_i^t-r_i^{t-1}\|_\infty^2,
 \qquad V_i(T)=\sum_{t=2}^TH(x_i^t,x_i^{t-1})^2.
\end{equation}
This inequality holds for arbitrary observed payoff vectors in $[0,1]^{m_i}$
and every fixed $\eta\leq1/16$, before the self-play coupling is imposed.

\begin{theorem}\label{thm:upper}
Suppose $|J_i|\leq D$ for every player, where $1\leq D\leq N$.
With $\eta=1/(24\sqrt D)$ the displayed rule satisfies
\begin{equation}\label{eq:sharpbound}
 R_i(T)\leq\frac{984}{13}\sqrt D(2+\log K)
             +\frac{15}{26\sqrt D}
 \leq76\sqrt D(2+\log K)
\end{equation}
for every player, game in this class, and horizon.
The rule uses $D$ but does not use the sets $S_i$ or the payoff tables.
\end{theorem}
\begin{proof}
Fix a finite $T$ throughout. At $t=1$, $\|r_i^1-r_i^0\|_\infty\leq1$.
Lemma~\ref{lem:variation} gives
\[
 \sum_{t=1}^T\|r_i^t-r_i^{t-1}\|_\infty^2
 \leq1+4\sum_{j\in J_i}V_j(T).
\]
Substituting in~\eqref{eq:rvu}, with $B=3c+10\eta^2$, yields
\begin{equation}\label{eq:local}
 \Psi_i(\eta U_i^{T+1})+\tfrac14 V_i(T)
 \leq B+40\eta^2\sum_{j\in J_i}V_j(T).
\end{equation}
Let $V_{\max}(T)=\max_iV_i(T)$. This is finite for finite $T$.
Choose a maximizing player in~\eqref{eq:local} and use positivity of
its potential. Then
\[
 \bigl(\tfrac14-40D\eta^2\bigr)V_{\max}(T)\leq B.
\]
Our rate has $\eta\leq1/24<1/16$ and $D\eta^2=1/576$, so the
coefficient is $13/72>0$. Returning to~\eqref{eq:local} for any player
and discarding only its nonnegative movement term gives
\[
 \Psi_i(\eta U_i^{T+1})
 \leq B\left(1+\frac{40D\eta^2}{1/4-40D\eta^2}\right)
 =\frac{18}{13}\left(3c+\frac{5}{288D}\right)
 =\frac{54}{13}c+\frac5{208D}.
\]
Since $R_i(T)=\max_aU_{ia}^{T+1}$, the potential lower bound gives
\[
 \eta R_i(T)\leq\Psi_i(\eta U_i^{T+1})-c
 \leq\frac{41}{13}c+\frac5{208D}.
\]
Division by $\eta$ proves the first inequality of~\eqref{eq:sharpbound}.
The difference $76-984/13=4/13$ and $c\geq2$, $D\geq1$ imply
\[
 \frac{15}{26\sqrt D}\leq\frac4{13}c\sqrt D,
\]
which proves the second. Every bound is uniform in the chosen finite
$T$; no infinite movement bound was assumed in the absorption step.
\end{proof}

\begin{corollary}
For the full problem,
\[
 1-\frac1K\leq C^*(N,K)\leq76\sqrt N(2+\log K).
\]
\end{corollary}
\begin{proof}
The upper bound follows with $D=N$, for every allowed action-count vector.
For the lower bound, fix an admissible tuple with some player having $K$
actions. Its first distribution $q$ is independent of the game.
Choose an action $a^*$ with $q(a^*)\leq1/K$ and give that player payoff
$1$ at $a^*$ and $0$ elsewhere, independently of the opponents.
Then $R_i(1)=1-q(a^*)\geq1-1/K$. Other players' payoffs may all be zero.
Take the required supremum and infimum in their original order.
\end{proof}

\section{An exact minimax subclass}
Call a game strategically separable if
\[
 u_i(a_i,a_{-i})=w_i(a_i)+h_i(a_{-i})
\]
for every player, with the total payoff in $[0,1]$. Neither function is
given to the learner. The functions $h_i$ may involve all opponents;
separability is a statement about payoff differences between own actions.

\begin{theorem}\label{thm:separable}
On the class of strategically separable games with fixed action vector
$m$, the minimax constant is $\max_i(1-1/m_i)$. Maximizing over
$m_i\leq K$ gives exactly $1-1/K$, for every number of players.
\end{theorem}
\begin{proof}
Use the uniform distribution at round one. The observed payoff vector
is $w_i+h_i(x_{-i}^1)\mathbf1$, so it identifies a maximizer $a_i^*$
of $w_i$. Break ties by the public action order and play that action
forever after. This rule is deterministic and uses only own feedback.
For any opponent sequence, all $h_i$ terms cancel from the regret.
The best fixed action is $a_i^*$, and its entire cumulative advantage is
\[
 \max_aw_i(a)-\frac1{m_i}\sum_aw_i(a).
\]
For any fixed opponent profile, normalized total payoffs imply
$\max_aw_i(a)-\min_aw_i(a)\leq1$. At least one summand is maximal,
so the displayed quantity is at most $1-1/m_i$.
The one-round lower-bound construction in the preceding corollary is
strategically separable, and applies to a player with largest $m_i$.
It matches the upper bound.
\end{proof}

\section{Padding and independent products}
\begin{proposition}\label{prop:monotone}
The function $C^*(N,K)$ is nondecreasing in $N$ and in $K$.
\end{proposition}
\begin{proof}
Increasing $K$ enlarges the finite set of action-count vectors in the
outer maximum; the admissible tuple for a fixed vector does not change.
To increase the player count from $N$ to $N'>N$, pad any $N$-player
game with $N'-N$ one-action players and make original payoffs independent
of them. Any tuple of $N'$-player rules then induces an admissible tuple
on the original $N$ players: each original learner runs its corresponding
rule with the padded public dimensions hard-coded. Its own payoff-vector
history is exactly the same. Dummy players have no effect on play or
regret. Thus the worst-case value for each padded action vector is at
least the value for its original vector. Take the outer maximum.
\end{proof}

\begin{proposition}\label{prop:products}
Fix a public partition of players into blocks $B_1,\ldots,B_s$ and a
public action-count vector. Consider only games in which a player's
payoff depends on actions in its own block. If $c_r$ is the minimax
constant for the unrestricted game on block $B_r$ with its prescribed
action counts, the minimax constant of this product class is
$\max_rc_r$.
\end{proposition}
\begin{proof}
For the upper bound, for any $\epsilon>0$ choose a block tuple with
worst-case regret at most $c_r+\epsilon$. Run the block tuples
simultaneously, translating labels using the public partition. A player's
feedback is exactly the feedback in the isolated block game. Consequently
the maximum individual regret is at most $\max_rc_r+\epsilon$.
Let $\epsilon$ decrease to zero.

For the lower bound, fix an arbitrary tuple for the product class and
one block $r$. Since a learner sees only its own feedback and internal
state, the restricted rules on that block form an admissible block tuple
with the remaining public data hard-coded. No other block can send it
feedback information. The supremum over games in that block is at least
$c_r$, by definition of $c_r$. Extend each such game by arbitrary fixed
games in the other blocks. This proves the lower bound $c_r$ for every
$r$, and therefore their maximum.
\end{proof}

\section{What the attempted solution does not establish}
The universal prefactor improvement and the sparse-game theorem leave
the asymptotic question open. In particular:
\begin{enumerate}[leftmargin=*]
\item No lower bound growing with $N$ or $K$ is established. The sharp
separable-game result cannot be used as an upper bound for coupled games.
\item Taking independent copies of a hard fixed-size game cannot amplify
the maximum individual-regret minimax value: Proposition~\ref{prop:products}
gives a maximum, not a sum.
\item The usual adversarial online-learning lower bound is not a lower
bound here without a reduction that realizes its payoff vectors in one
fixed game while all other players follow their prescribed rules.
No such reduction is supplied.
\item The sparse bound uses a rule tuned to the promised class. It does
not establish adaptation to an unknown interaction bound while retaining
the same guarantee in all games.
\end{enumerate}
A resolution of WS-107790 still requires a matching dimension-dependent
lower bound or an improvement of the unrestricted asymptotic upper bound.

\section*{Audit and provenance}
The numerical absorption, first-round term, $T=1$ case, one-action
players, and quantifier order were checked. All new mathematical claims
are proved above, apart from the explicitly cited MORM estimate.
The quoted $96\sqrt N(2+\log K)$ result is credited to its authors;
this contribution does not claim discovery of constant-regret learning.
\section*{Models used}
OpenAI GPT 6 Astra Pro generated the attempt and performed proof checks,
including separate-agent review of the potential absorption argument.
Exact model build and reasoning setting are unknown.

\section*{Human contributors and roles}
The user selected the problems, supplied the package, and requested the
attempts. No human verification or external referee approval is claimed.
The literature results retain the authorship indicated below.

\begin{thebibliography}{9}
\raggedright
\bibitem{problem}
We-Solve. \emph{Optimal dimension dependence of constant self-play regret}.
Problem supplied as WS-107790; live intake handle ECON-1. Statement and
research prompt read 18 October 2026.
\url{https://we-solve.github.io/problems/optimal-dimension-dependence-of-constant-self-play-regret/}.
\bibitem{morm}
Ashkan Soleymani and Georgios Piliouras.
\emph{Multiplicative Optimism for Constant Regret in Games}.
arXiv:2609.21976v1, 18 September 2026. Used: Algorithm 1,
Lemma 4.1, Proposition 4.7, and Theorem 4.10. The HTML and PDF were
retrieved and the cited estimates inspected.
\url{https://arxiv.org/html/2609.21976v1}.
\end{thebibliography}
\end{document}
